\section*{Chapter \ref{ch:g7:microscope} --- The Microscope and the Cell}

\begin{solution}{exo:g7:microscope:1}
A \omterm{def:g7:microscope:cell}{cell} is the smallest unit of an \omterm{def:g4:ecosystems:levels}{organism} that can carry out basic life
activities. \omterm{def:g7:microscope:cell}{Cells} are usually tens of micrometres or smaller --- below
comfortable unaided resolution --- so a \omterm{prop:g4:cells:building}{microscope} is needed to see them
clearly.
\end{solution}

\begin{solution}{exo:g7:microscope:2}
Eyepiece: further magnifies the image for the \omterm{def:g5:nerves:senses}{eye}. Objective: first lens
near the specimen; main \omterm{def:g7:microscope:mag}{magnification}. Stage: holds the slide. Coarse
focus: large focus adjustments (low power). Fine focus: small sharp
adjustments. Diaphragm: \omterm{def:g6:method:control}{controls} \omterm{def:g4:plants-need:needs}{light}/contrast.
\end{solution}

\begin{solution}{exo:g7:microscope:3}
Total \omterm{def:g7:microscope:mag}{magnification} $= 10 \times 40 = 400$.
\end{solution}

\begin{solution}{exo:g7:microscope:4}
Both: \omterm{def:g7:microscope:parts}{cell membrane}, \omterm{def:g7:microscope:parts}{cytoplasm}, \omterm{def:g7:microscope:parts}{nucleus} (and ribosomes, mitochondria ---
if known). Plant-only (typical): \omterm{def:g7:microscope:parts}{cell wall}, \omterm{def:g7:microscope:parts}{chloroplasts}, large central
\omterm{def:g7:microscope:parts}{vacuole}.
\end{solution}

\begin{solution}{exo:g7:microscope:5}
(1) \omterm{def:g4:ecosystems:levels}{Organisms} are made of \omterm{def:g7:microscope:cell}{cells}; (2) \omterm{def:g7:microscope:cell}{cell} is basic unit of structure and
function; (3) \omterm{def:g7:microscope:cell}{cells} come from \omterm{def:g7:microscope:cell}{cells} by division.
\end{solution}

